\chapter{Aide et accompagnement}

Au cours de votre utilisation du logiciel Tockem SPE, il est possible que vous rencontriez des difficultés ou des obstacles susceptibles d'entraver votre progression.

Dans un premier temps, nous vous recommandons vivement de consulter à nouveau ce document, qui regroupe l'ensemble des informations nécessaires à l'utilisation du logiciel. La solution à votre problème s'y trouve très probablement, expliquée de manière claire et détaillée.

Si, malgré cette relecture, vous ne parvenez toujours pas à résoudre votre souci, nous vous invitons à contacter l'équipe dédiée du volet Eau, Hygiène et Assainissement (EHA) de l'association TOCKEM. Cette équipe, composée d'experts formés et engagés, est à votre disposition pour analyser votre situation et vous fournir une solution adaptée dans les meilleurs délais.

De la même manière, si vous souhaitez proposer l'ajout de nouvelles fonctionnalités au logiciel Tockem SPE afin d'optimiser davantage la gestion des réseaux d'adduction d'eau potable, n'hésitez pas à soumettre vos suggestions à l'équipe EHA. Vos retours et idées sont précieux pour faire évoluer cet outil et répondre au mieux aux besoins spécifiques de la commune de Fokoué, contribuant ainsi à une gestion plus efficace et durable des ressources en eau.
